\begin{abstract}

This thesis investigates whether a mechanical deflection mechanism---a freely rotating protective cage---can reduce impact shock and vibration in confined-space drone operations, and whether the onboard inertial signatures of those collisions can be exploited for spatial awareness.

A custom 4-inch autonomous quadcopter was designed and built, integrating a Pixhawk~6C flight controller running PX4, a Raspberry~Pi~5 companion computer running ROS~2, and a PETG spherical cage (35.8\,cm diameter) tested in two configurations: rigidly fixed and bearing-decoupled (freely rotating). All other parameters---mass, diameter, material, and propeller clearance---were held constant, isolating the bearing mechanism as the single independent variable. Experiments were conducted in an OptiTrack motion-capture environment, with the drone executing 179 autonomous collision passes against a static cylindrical obstacle at approach angles of $45^\circ$ and $75^\circ$. After telemetry-based verification, 127 structural impacts were isolated for comparative analysis.

The rotating cage reduced maximum impact deceleration by 59\% relative to the fixed configuration ($-1.64 \pm 0.64$\,m/s$^2$ vs.\ $-1.03 \pm 0.50$\,m/s$^2$, $p < 0.001$, Cohen's $d = 1.05$), and IMU vibration analysis revealed a narrower post-impact oscillation envelope with faster return to the nominal noise floor. The bearing mechanism introduced a modest spatial trade-off: maximum trajectory deviation increased slightly (Fixed: $118.4 \pm 30.4$\,mm; Rotating: $131.0 \pm 32.1$\,mm, $p = 0.022$, $d = 0.41$), but the total recovery area remained statistically indistinguishable ($p = 0.998$), indicating that the flight controller's PID correction capacity absorbed the offset.

To evaluate whether impact angle can be inferred from onboard sensors alone, a Random Forest regressor was trained on 14 IMU-derived features. Cross-validated $R^2$ reached approximately 0.72 for both configurations, with mean absolute error of $4.5^\circ$--$4.9^\circ$---sufficient to distinguish glancing from deeper collisions. Critically, the rotating cage's angle signature was invisible to linear models (Huber baseline $R^2 = 0.004$) yet fully recovered by the non-linear Random Forest ($R^2 = 0.719$), and cross-condition transfer between cage types failed decisively ($R^2 < 0.26$), demonstrating that the bearing mechanism fundamentally redistributes---rather than destroys---the angle information carried by the vibration signature.

These findings establish the rotating cage as an effective mechanical countermeasure for impact survivability and open the door to \emph{tactile SLAM}: navigation by contact in environments where no other sensor modality is viable.

\end{abstract}
